\chapter*{Conclusions of Part~I}
\addcontentsline{toc}{chapter}{Conclusions of Part~I}
\label{ch:part1-concl}

The six chapters in this Part have addressed both the theoretical foundations and the observational consequences of WI, from identifying the necessary conditions for a viable WI scenario and developing the tools to confront it with CMB data, to uncovering novel signatures with no analogue in CI.

At the foundational level, we showed in Chapter~\ref{ch:wi-scalar-field} that the fine-tuning requirements on the inflaton potential are generally more stringent in WI than in CI, particularly for dissipation rates with non-negative temperature dependence ($\Upsilon \propto T^c$ with $c \geq 0$). While the enhanced friction reduces the required field excursion, the thermal amplification of scalar perturbations constrains the height of the potential more severely, resulting in bounds on the fine-tuning parameter $\lambda$ that are at least three orders of magnitude tighter than in the cold case. This counter-intuitive result clarifies that the additional friction in WI reduces the required field excursion but does not relax the constraints on the steepness of the potential, as previously indicated in the literature.  It also reinforces the appeal of axion-like inflaton constructions discussed in the introduction,  in which a shift symmetry protects the potential from the radiative corrections that would otherwise drive its self-coupling to large values, resolving the fine-tuning problem at the level of the underlying microphysics.

A related theoretical question is whether WI can avoid the onset of eternal inflation. In Chapter~\ref{ch:wi-eternal}, I derived a generalized bound on the running of the spectral index $\alpha_s$ required to prevent the emergence of a fluctuation-dominated regime. In WI the threshold for eternal inflation is raised relative to CI, meaning that an even smaller magnitude of negative running suffices to prevent it. I demonstrated that several physically motivated WI constructions naturally produce a running $\alpha_s\lesssim -10^{-4}$ to $-10^{-3}$, comfortably satisfying the derived bound while remaining consistent with CMB observations.

Confronting these theoretical predictions with data requires accurate numerical tools. In Chapter~\ref{ch:warmspy}, I developed the public code \texttt{WarmSPy} for computing the full primordial power spectrum in WI, enabling a systematic numerical study of cosmological perturbations across the full range of dissipative regimes. A central result is the universality of the scalar dissipation function $G(Q)$, which parametrizes the enhancement of the scalar power spectrum relative to CI, and which depends only on the dissipation strength $Q$ and the temperature power-law index $c$ --- independently of the inflaton potential, the field-dependence of the dissipation rate, and the inflationary history. The fitting formulae provided extend to $Q \sim 10^4$, covering the full range relevant to current WI constructions.

With these tools in hand, I applied the WI framework to natural inflation in Chapter~\ref{ch:wi-natural}. The original cosine potential, now in significant tension with CMB data in the CI context, can be brought into agreement with observations in the WI scenario, with the axion decay constant reaching marginally sub-Planckian values ($f \gtrsim 0.3\text{--}0.8\,M_\mathrm{pl}$ depending on the dissipation model) but not significantly below the Planck scale. Increasing the dissipation strength reduces the field excursion but simultaneously shifts the spectral index towards bluer values, tightening the viable parameter space.


Beyond refining the theoretical framework, the final two chapters explored novel phenomenological consequences that arise uniquely from the presence of a thermal bath during inflation. In Chapter~\ref{ch:wifi-dm}, I showed that this persistent radiation environment provides a natural and efficient setting for dark matter production via ultraviolet freeze-in, a mechanism we termed WIFI.\footnote{Full credit for this name goes to my advisor, Katherine Freese, whose instinct for branding is always sharp}. The resulting dark matter yield is always enhanced compared to production during a standard radiation-dominated era at the same reheating temperature, with the enhancement growing dramatically with the mass dimension of the non-renormalizable operator mediating the interaction. For higher-dimensional operators essentially all the dark matter is produced during the inflationary phase itself. In Chapter~\ref{ch:thermal-gravitons}, I demonstrated that particle collisions in the WI thermal bath also inevitably produce thermal gravitons, giving rise to a stochastic gravitational-wave background with a distinctive spectral shape consisting of a sharp peak near $\sim 100\,$GHz from sub-horizon thermal emission and a nearly flat plateau extending to low frequencies from super-horizon modes --- the latter being entirely absent in CI. Both the overall energy density and the low-frequency plateau are enhanced by one to two orders of magnitude compared to production during a standard radiation-dominated era; nevertheless, the predicted frequencies and amplitudes make detection with current and near-future experiments a significant observational challenge.

Taken together, these results highlight the rich phenomenology of WI and the new observational channels it opens into the inflationary epoch --- from dark matter production via freeze-in to a distinctive thermal gravitational-wave background. Beyond the signatures explored in this Part, WI is also expected to produce a sizable level of non-Gaussianity with a spectral shape that differs qualitatively from that generated in a CI setting~\cite{Bastero-Gil:2014raa, Mirbabayi:2022cbt}. Sharpening these predictions is particularly timely in light of the improved sensitivity expected from ongoing and upcoming space missions such as SPHEREx~\cite{SPHEREx:2014bgr} and Euclid~\cite{Euclid:2024ris}, as well as from CMB measurements by the Simons Observatory~\cite{SimonsObservatory:2018koc,SimonsObservatory:2020noise}, which will simultaneously tighten constraints on the tensor-to-scalar ratio and the spectral tilt across the viable WI parameter space.